\documentclass[11pt]{article}
\usepackage[a4paper,margin=22mm]{geometry}
\usepackage{fontspec}
\setmainfont{LinLibertine_R.otf}[BoldFont=LinLibertine_RB.otf,ItalicFont=LinLibertine_RI.otf,BoldItalicFont=LinLibertine_RBI.otf]
\usepackage{xcolor}
\usepackage{hyperref}
\hypersetup{colorlinks=true,urlcolor=blue}
\setlength{\parindent}{0pt}
\setlength{\parskip}{6pt}
\sloppy
\emergencystretch=2em
\begin{document}
\section*{Historical Britannica: extracted source}
Encyclopaedia Britannica - Eighth edition, Volume 3, Anatomy-Astronomy - EB.16\par
{\small Source row 744. Parser: nvidia/nemotron-3.5-parse.\par}
{\small This is model-extracted text, not a corrected transcription. OCR errors, reading order and region alignment require human review.\par}
\hrule\medskip
{\small\itshape 786 VASIA.\par}

{\scriptsize\color{gray} Region region-0; provider order 0.\par}
\medskip

Ata.
\par

{\scriptsize\color{gray} Region region-1; provider order 1.\par}
\medskip

the best known. Samuton represents the country as fertile and well cultivated, the people a industrious and critical. They, as well as the Tunkness are of Chinese descent. Their coasts abound with the edible leaf nests, formed by a species of swallow. See Coornis-Gunka and Trexus.
\par

{\scriptsize\color{gray} Region region-2; provider order 2.\par}
\medskip

China.
\par

{\scriptsize\color{gray} Region region-3; provider order 3.\par}
\medskip

The vast majority of China presents an area computed at 129/509, almost 20 years old among other countries, according to Sanuman, as a population of 333,000, being the most densely populated region of the earth. Its enormous surface is much varied, pro- during the vegetable richer of every climate. Chinese Tatarty is very common), and claims of granite monuments traverse China in various directions; but the great chains generally near from W. to that said such the principal rivers in that direction through fertile plains of enormes extent. The principal of these trees are the Juan, in the north; the Hanhoky, and Kian-ku in the central districts, both having a coarse of more than 5000 miles; and the Han-kang, all and Ye-kan in the X. It is needless to state here what is discussed fully under Onyx.
\par

{\scriptsize\color{gray} Region region-4; provider order 4.\par}
\medskip

Oceania Asia endowed all the Natali Bare properly so called, and some would also include under this hand the distant New Zealand, and vast continent of Australia. We limit this division of Greyion, the Ishaben Ides, the Son matran chain, including Java, Sec., Borneo, Lobbers, the Moscow or Spice Islands, the Muslims of Palinyans, Par- mosu, Luk-kna, Japan, sees, and Il-Toma. For a particular account of this division, we refer to the different heads just enumerated.
\par

{\scriptsize\color{gray} Region region-5; provider order 5.\par}
\medskip

{\small\itshape Climate\par}

{\scriptsize\color{gray} Region region-6; provider order 6.\par}
\medskip

Asia, extending from the equator to the Aretie Sea, necessarily possesses great variety of climatic. But though here as its over other country, the climate is regulated by the distance from the equator, this general law is modified by accidental causes. In the wide extent of Africa versus recent production on temperate cover, which cannot be very clearly explained. To such inquiry some uncertainty will always attack, and anomalies may appear of which we can offer no solution. Tests are not only sure guides it and unless thereby, in the following observations, we will endeavour to address.
\par

{\scriptsize\color{gray} Region region-7; provider order 7.\par}
\medskip

The height of the land above the level of the sea is an sure it cause of cold as distance from the equator; and counter one not only cold in proportion to their height, but a mass of cold air is accumulated above them, which, being dispersed, is carried towards the equator, or extends the dominion of cold into the regions of heat. Land and water, also, or the causes either of heat or cold, according to their situation. The green areas of the oceans is little affected by the changes of the seasons; and is consequently preserves the medium temperature of the whole year. Hence the vicinity of the ocean ends the temperature of the equinoctial regions and in higher latitudes it underlies the extreme both of heat and cold, being the winter of a higher temperature, and in summer of a lower temperatures, than the momentumbratir. The two surface of the earth, again, imbishes heat or cold, much more readily than the ocean; and it is only a considerable depth that is found to give the median temperature of the year. The vicinity of land, therefore, in the polar countries, is the cause of cold, while in the southern regions where the equator it is an equally power-ful cause of heat. Thus Africa, which extends so far to the S, and which contains a greater proportion of land within
\par

{\scriptsize\color{gray} Region region-8; provider order 8.\par}
\medskip

the tropics than any other division of the globe, is a vast stone-house of heat, from which it is dispersed for and wide, and even reaches the shores of Europe in hot and parching winds; while, on the other hand, the beach and extended the Anamon continues toward the V sufficiently accounts for the coldness of its climate--the N.W. winds which sweep across its frozen wastes exceeding their rounds into the regions of heat as its sometimes as Mexico or Yera Cruz. The influence of the ocean in modeling the severity of the winter is completed in the climate of Great Britain, where no such intense cold ever prevails as in corresponding latitudinales on the continent of Europe, and more especially on the of Asia. The climate of country has affected by the direction of the winds; and hence the eastern shores of Antarctica, owing to the trade winds, which blow from the E, and are cooled in their passage across the climatic, have not the same starkly basised the opposite shore of Africa, where the same roads are heated to the intense degrees in their pursage over the burning deserts of the interior. The northern frontiers of Asia, and its prolificable elevation towards the two, necessarily enough the greatest portion of its to the dominion of cold. Among the crustal structures perpetrity and regime, and from these snowny deserts the influence of cold is widely extended over the high plains of the intor- rior. This they wish is about the same latitude as Northern Africa of Arabia, namely between the 2001, and 25th de- greers, there is a continuous and severe winter of three months, which is of each uniforms recently that at its commencement the inhabitants till their most, and it is kept perfectly fresh for the moments." To the W, along the white map of ele- tentooman that extends into Persian and the Caspian, the climate is modified by the elevation of the ground. In the countries of Balk and Pakistan, should lie on the western delsing of the great reign of the Hindu Koh Mountains, in the same latitude as the S of Europe, many, the 20th and 40th degrees, and all along the banks of the Ova, the climate is remarkably severe. For three months the winter is intensely cold, the wind being dry and affecting and the snow lying close on the ground. The rivers are all frozen over, and the Oxus during all that periods possible for ear- vans." The summer, again, is exactly her. Perks, in like manner, being nearly in the same latitude as Arabia, the hottest part of the earth, and lasting as excessively hot some lies, last the northern and central parts of several winter of a northern climate, with drifting snow, which lies deep on the ground for the few months; and this is largely entirely to its elevation, which is estimated by France to be 4000 feet above the Caspian Sea." The lower valleys and exterior plains of Asia, which lies to the west of the Himalaya Morou tains, including Arabia, the southern and flat parts of Per- as, Thailand, and India beyond theanges, considering the tropical and warm regions of this continent, of which the climate, though it agrees to ground with their position on the globe, still varies from local causes. Hindulus, for example, and India beyond the Ganges, though they approach nature to the equator, nor we now rely our hot as Arabia or the glaciers countries. The course of the economic is also more constant; and it is here that we meet with those remarkable winds, the mountains, and hills here as months in opposite directions, from the N.W. and N.V., with some slight variations, and which extend their influence over all the countries
\par

{\scriptsize\color{gray} Region region-9; provider order 9.\par}
\medskip

3 See Humbel's Political Essory in the Kingdom of New Spain. He mentions that at Yera Cruz: the emigrated thermometer som- times falls, owing to those northern blasts, to 0 or to 52 degrees of Palaresheit; and that even at Mexico they sometimes, though rarely, bring down the thermometers to the freezing point.
\par

{\scriptsize\color{gray} Region region-10; provider order 10.\par}
\medskip

2 Towards Advisory to Turkey, p 51 / 591, 592.
\par

{\scriptsize\color{gray} Region region-11; provider order 11.\par}
\medskip

3 Presently Secretary of Joanneya and Alexander, Aprendiz, p . 50.
\par

{\scriptsize\color{gray} Region region-12; provider order 12.\par}
\medskip

4 Kimin's Appropriekal Minister of the Vienna Dayire, p 122. It is mentioned in the Journal of a Mr Campbell, quoted by Sir J. Malcolm, that while he was at Talverez a heavy fall of snow occurred in May; and in December and January the thermometer at night was never above zero. (Sir J. Malcolm's History of Texteio, vol. IL p. 000.)
\par

{\scriptsize\color{gray} Region region-13; provider order 13.\par}
\medskip
\newpage\section*{Source record}
bri-0744-fcdb217a9ba5\par
Image SHA-256: \texttt{fcdb217a9ba5e3b917823cdc39f41c29}\par\texttt{fcfee28beaacefc676ceeef419d1b70f}\par
Dataset revision: \texttt{309f765b1ecc83e0a7f7}\texttt{762e44b4a927a7555443}\par
National Library of Scotland, Encyclopaedia Britannica Illustrated (CC0).\par
\url{https://huggingface.co/datasets/NationalLibraryOfScotland/encyclopaedia_britannica_illustrated/viewer/default/train?row=744}
\end{document}
